\chapter{Dependencies, Tests, Examples, and Developer Tools}
\label{ch:dependencies-tools}

A successful CNA configure can involve checked-out submodules, sibling repositories, system
packages, downloaded source, prebuilt binary releases, and renderer-specific SDKs.  Treating all of
them as “vendored dependencies” hides the most common reproducibility failures.  This chapter maps
the acquisition mechanisms and the developer surfaces that consume them; Chapter~\ref{ch:configuring-building}
owns the command-line configuration details.

\section{Five acquisition mechanisms}

\begin{center}
\small
\begin{tabularx}{\textwidth}{lXX}
\toprule
Mechanism & Representative inputs & Reproducibility boundary \\
\midrule
Git submodules & GoogleTest, SDL3, SDL3\_image, SDL3\_mixer & A source archive may omit them;
  initialize the recorded submodule commits before configuring. \\
In-tree sources & cgltf, stb, ENet, and a wgpu-native vendor area & Present in the CNA tree, but
  build/features still depend on the selected module and host. \\
Sibling checkouts & sharp-runtime, easy-gl/meta-gl, free-direct/free-api & Separate repositories
  with their own pins and sometimes fixed sibling-relative paths. \\
Fetch/download or external artifact & bgfx, Blend2D, Diligent, FNA3D, LLGL, Magnum, OpenVG,
  PortableGL, sokol, WebGPU, Wicked, Skia & Each renderer has a distinct pin, offline override,
  patch, or prebuilt-root contract. \\
System packages & FFmpeg, optional Draco, Vulkan/OpenGL/GLES libraries, Threads & Host package
  discovery can change which files compile or whether configure is admissible. \\
\bottomrule
\end{tabularx}
\end{center}

The categories can compose.  FNA3D is fetched from a pin and carries MojoShader through its own
submodules.  EasyGL is a sibling checkout that itself expects meta-gl as a sibling.  Free-direct is
a sibling that PUBLIC-links the second-level free-api sibling.  A lockfile for only CNA's own git
tree is therefore not a complete build identity.

\section{Submodules: initialize, do not guess}

CNA's considered instruction is \texttt{git submodule update --init}, not a recursive update.
The optional codec trees beneath SDL are disabled by CNA's vendored-SDL configuration, so recursive
initialization adds time without contributing compiled inputs.  The configure guards explain when
the four required submodule directories are missing and explicitly note that a plain source archive
cannot manufacture their content.

For a distributable source snapshot, either include those exact submodule trees or select a
supported system-SDL path and provide GoogleTest another documented way.  Silently building against
whatever SDL happens to be installed changes the dependency claim.

\section{The unusual SDL configure-time build}

The controlling option is \texttt{CNA\_USE\_SYSTEM\_SDL}.  With the default vendored path, CMake
configures, builds, and installs SDL3, SDL3\_image, and SDL3\_mixer during CNA's
\emph{configure} step.  The install prefix is persistent
and outside the individual build directory, then consumed through the packages' exported CMake
configs.

This has three operational consequences:

\begin{itemize}
  \item deleting or cleaning a CNA build directory does not necessarily rebuild SDL;
  \item the prefix is keyed by target system and processor so a cross-build does not overwrite a
        native installation;
  \item a configure failure may actually be a nested dependency compile/install failure, not a
        CNA CMake-language error.
\end{itemize}

The default parallelism ceiling is intentionally small.  Image and mixer codec options are pruned:
for example, the vendored SDL\_mixer build does not enable the optional MP3, Opus, FluidSynth,
WavPack, or FLAC dependency families.  CNA's higher-level Audio and Media behavior must be read
against that concrete mixer build, not against every feature SDL\_mixer can support in principle.

\section{System-package gates can change source composition}

FFmpeg illustrates why “dependency found” is not merely a link flag.  CNA computes the
\texttt{CNA\_FFMPEG\_AVAILABLE} flag.  On Windows/MinGW, Emscripten, and Android it is forced off
so a
cross-configure cannot accidentally import the host's native pkg-config headers and libraries.
When off, the Media module drops its video decoder/player translation units, Content drops the XNB
video reader, and the unit-test source set excludes matching cases.

Draco is optional in a different way.  Its absence does not remove the glTF importer; a
Draco-compressed primitive reaches a named unsupported-format failure.  Vulkan and several GL-based
renderers, by contrast, require their native SDK/package at configure time.  These are three
different states: feature source omitted, format rejected at runtime, and configuration refused.

\section{Sibling repositories are part of the workspace contract}

sharp-runtime is always required.  \texttt{CNA\_SHARP\_RUNTIME\_ROOT} defaults to
\texttt{../sharp-runtime}, and CNA emits a fatal diagnostic explaining that it is not a submodule.
The consumption adapter detects the current component-target shape and links each CNA module to a
declared sharp-runtime component closure; it can also recognize the older monolithic target shape.
For this edition, only the revisions in Appendix~\ref{ch:repo-map} are authoritative.

EasyGL's five public identities expect \texttt{../easy-gl}; easy-gl then expects
\texttt{../meta-gl}.  CNA exposes no override for those two relative paths.  FREEDIRECT similarly
expects the free-direct sibling checkout, which resolves free-api itself and reuses CNA's
already-created SDL
targets.  A CI checkout that clones only CNA can therefore configure some renderer identities but
not others.

Do not paper over this by copying sibling sources into CNA's build tree.  Their independent commit,
tests, public targets, and consumer evidence are part of the claim.

\section{Renderer dependencies are intentionally asymmetric}

Some renderer integrations fetch a pinned source tree automatically; some prefer a supplied root;
some require a prebuilt external product.  Examples at the pin include:

\begin{itemize}
  \item WEBGPU auto-downloads a pinned wgpu-native binary release unless an explicit root is given;
  \item WICKED has a pinned commit and patch set but auto-fetch is off by default;
  \item SKIA is never fetched and requires caller-supplied prebuilt directories;
  \item MAGNUM first tries system Corrade/Magnum packages, then falls back to pinned source;
  \item BGFX applies an in-tree patch during its fetched-source preparation;
  \item OPENVG and several Sokol modes still require a real OpenGL system package.
\end{itemize}

Accordingly, “all renderer identities accepted by CMake” is not “all renderer dependencies are
available on this host.”  Chapter~\ref{ch:renderer-selection-identity} separates identity,
family, internal native sub-selector, and platform admission.

\section{What CNA builds, and what it does not package}

CNA provides build-tree targets.  At the pinned revision it has no install rules, exported package
configuration, or supported \texttt{find\_package(CNA)} consumer flow.  A parent project normally
adds the checkout as a subdirectory or follows a sibling template that does so.  Inventing an
installation step produces an untested include/library layout and loses the sibling-resolution
contract.

Likewise, \texttt{CNA} is a CMake target, while \texttt{CNA::Math} and other namespaced aliases are
module targets.  Neither spelling implies a file named \texttt{libCNA.so}; most framework pieces are
static archives composed into the final consumer.

\section{The CnaTests corpus}

The main GoogleTest executable is assembled from module-owned test sources plus the selected
renderer configuration.  Feature gates, platform exclusions, renderer selection, and optional
FFmpeg/Net inputs change its compiled definition population.  There is no configuration-free
“number of CNA tests.”  Chapter~\ref{ch:cnatests-architecture} explains composition, and
Chapter~\ref{ch:counting-discipline} derives the pinned static populations without exchanging files,
macros, CTest registrations, and executed cases.

The native \texttt{tests} preset deliberately says to run \texttt{CnaTests} directly.  CTest's
GoogleTest discovery can launch each case as a short-lived process, while this repository includes
display-dependent binaries and fixtures sharing fixed temporary paths.  CTest remains useful for
named standalone gates and renderer executables; it is not automatically the authoritative runner
for the monolithic corpus.

\section{Examples are executable documentation, not one evidence tier}

Module-local examples range from tutorial programs to GPU readback tests, golden-image comparisons,
host/browser harnesses, and policy-only probes.  The filename or registration directory does not
tell which one it is.  For each example record:

\begin{itemize}
  \item whether it is compiled by the selected configuration;
  \item whether it runs automatically or only by explicit dispatch;
  \item which renderer/native driver/host actually engaged;
  \item whether its oracle is exit status, structured values, pixels, or human inspection;
  \item which unavailable-host outcome is skip, refusal, or failure.
\end{itemize}

Shared golden images remain at repository level because many renderer modules consume the same
oracle.  A renderer-owned test executable plus a shared image is still two artifacts whose
provenance and update rules must travel together.

\section{Presets are named environments, not universal promises}

The pinned presets cover a Web build, three Devices sanitizer builds, and a native Ninja test
build.  The sanitizer presets pass raw compiler/linker flags even though CNA also exposes the
\texttt{CNA\_SANITIZE} option; this is redundant but observable.  They explicitly enable
\texttt{CNA\_DEVICES}, preventing a sanitizer job named for Devices from silently compiling the
feature out.

A preset proves only its declared target and commands.  The Web preset builds the house demo with
tests off; it is not a browser conformance run.  The native test preset builds \texttt{CnaTests}; it
does not make every standalone CTest executable.  Hosted workflow reality is audited separately in
Chapter~\ref{ch:ci-reality}.

\section{Developer tools and what they decide}

CNA's mechanical tools fall into distinct authority classes:

\begin{description}
  \item[Structural validators] Compare finite source-owned sets, paths, and identity
        spellings.  They are strong for structural equality and silent about behavior.
  \item[Module probes and closure checkers] Compile a minimal consumer or inspect a link artifact.
        Their verdict depends on generator and host engagement, as Chapter~\ref{ch:boundaries-enforcement}
        shows.
  \item[Content and catalog checkers] Validate CNJ/glTF fixtures, examples, or documented catalogs.
        Their parser and denominator define the result.
  \item[Renderer/golden harnesses] Exercise a chosen native path and compare pixels or structured
        output.  They require identity and driver evidence in addition to a green comparison.
  \item[Cross/toolchain launchers] Stage DLLs, select Wine prefixes, start browsers, or prepare
        target files.  Successful process creation is not yet proof that the intended runtime did
        the work.
\end{description}

The durable rule is to keep acquisition, build, registration, execution, and oracle engagement as
separate facts.  A dependency pin enables reproducibility; it does not establish that a test used
the dependency.  A built example establishes reachability; it does not establish correct pixels.
Appendix~\ref{ch:verification-tiers} provides the compact claim vocabulary used throughout the rest
of the book.
